\paragraph{Cognitive models of writing.}
Writing-process research asks how planning, text production, and evaluation are coordinated during composition. The cognitive process theory of writing~\cite{flower1981cognitive} is particularly useful for our purpose because it assigns this coordination to a \texttt{Monitor} and allows writer-generated goals to change as writing proceeds. Later accounts distinguish knowledge-telling from knowledge-transforming~\cite{bereiter1987psychology}, model text production as knowledge generation~\cite{galbraith1999writing}, and revise the original architecture~\cite{hayes1996new,hayes2012modeling}. We treat these accounts as theories of writing organization from which computational hypotheses can be derived, not as descriptions of a language model's internal cognition.

\paragraph{Structured long-form generation.}
Long-form generation systems also make decisions beyond token-level generation, but usually about \emph{what content or task to work on next}. Plan-and-Write and Re$^3$ use plans to organize subsequent realization~\cite{yao2019plan,yang2022re3}; STORM builds an outline before section-level drafting~\cite{shao2024assisting}; LongWriter's AgentWrite decomposes a long response into manageable subtasks~\cite{bai2025longwriter}; and WriteHERE recursively expands and executes an evolving task graph~\cite{xiong2025beyond}. Self-Refine repeatedly applies feedback and revision~\cite{madaan2023self}, while IS-CoT cycles through planning, writing, and reflection~\cite{sun2026cot}. These systems vary substantially in scheduling and feedback, but their next-step decisions are typically tied to a plan element, section, passage, task, or prescribed stage. Our formulation instead makes the choice among writing processes itself the controller's action space.

\paragraph{Process-oriented and cognitively inspired writing systems.}
Writing-process theory has also informed systems for human--AI co-writing and writing support~\cite{gero2022design,wan2024coco,siddiqui2025script,zhou2026iris}. CogWriter~\cite{wan2025cognitive} is the closest autonomous system: it combines hierarchical planning, generation, monitoring, and reviewing under an explicit cognitive-writing motivation. Our distinction is therefore not cognitive inspiration or the presence of planning and review. It is that the controller explicitly decides \emph{which writing process to execute next}, and that the commitments defining this decision---explicit goals, adaptive process ordering, and separated process roles---are exposed as interventions. Appendix~\ref{app:mechanism-comparison} compares these mechanisms with representative systems. This perspective also connects to language-agent work on explicit control loops and role-structured computation~\cite{sumers2024cognitive,yao2023react,shinn2023reflexion}.
